\scanpage{9}
\begin{center}Sur les groupes de Galois attachés aux groupes $p$-divisibles\hfill123\end{center}

\begin{center}\textbf{\S\ 3. Le groupe de Mumford-Tate $p$-adique}\end{center}

Les notations étant celles du \S\ 2, soit $V$ un $\mathbf Q_p$-espace vectoriel de dimension finie muni d'une loi d'opération de $\operatorname{Gal}_K$ continue et linéaire. Soit $G$ l'image de $\operatorname{Gal}_K$ dans $\operatorname{Aut}(V)$; c'est un sous-groupe compact de $\operatorname{Aut}(V)$.

Faisons les deux hypothèses suivantes:

(H.1) -- \textit{$V$ est un $G$-module semi-simple.}

(H.2) -- \textit{Le module galoisien $V_C=V\otimes_{\mathbf Q_p}C$ est du type Hodge-Tate.}

(Précisons que $\operatorname{Gal}_K$ opère sur $V_C$ par transport de structure, i.e. par la formule $s(v\otimes c)=s(v)\otimes s(c)$.)

L'hypothèse (H.1) permet d'appliquer à $G$ les définitions et résultats du \S\ 1; en particulier, le groupe $G_{\mathrm{alg}}$ est défini. On notera que $G_{\mathrm{alg}}(C)$ est un sous-groupe du groupe $\operatorname{Aut}(V_C)$ des automorphismes $C$-linéaires de $V_C$.

Soit d'autre part $\lambda\in C^*$; notons $\varphi(\lambda)$ l'automorphisme de $V_C$ qui est l'homothétie de rapport $\lambda^i$ sur la $i$-ème composante $V_C(i)$ de $V_C$ (au sens du \S\ 2). On définit ainsi un homomorphisme $\varphi:C^*\to\operatorname{Aut}(V_C)$. Soit $\Phi=\operatorname{Im}(\varphi)$; c'est un sous-groupe de $\operatorname{Aut}(V_C)$.

\textbf{Théorème 1.} -- \textit{Le groupe $\Phi$ est contenu dans $G_{\mathrm{alg}}(C)$.}

Soit $(r,s)$ un couple d'entiers $\geq0$, et soit $W=T_{r,s}(V)$. Vu la prop. 3, il suffit de prouver que tout élément de $W$ invariant par $\operatorname{Gal}_K$ est invariant par $\Phi$ (lorsqu'on l'identifie à un élément de $T_{r,s}(V_C)=W_C$). Or:

(a) \textit{$W_C$ est du type Hodge-Tate.}

(b) \textit{Si $\lambda\in C^*$, l'action de $\varphi(\lambda)$ sur $W_C$ est donnée par:}
\[
\varphi(\lambda)x=\lambda^i x\quad\textit{si}\quad x\in W_C(i).
\]
En effet, (a) et (b) sont vrais pour $W=V$ (i.e. $r=1,s=0$), et, si elles sont vraies pour $W_1$ et $W_2$, elles le sont aussi pour $W_1^*$ et $W_1\otimes W_2$, en vertu de la remarque de la fin du \S\ 2.

Ceci étant, si $w\in W$ est invariant par $\operatorname{Gal}_K$, $w$ appartient \textit{a fortiori} à $W_C^0$, donc aussi à $W_C(0)=W_C^0\otimes_K C$, et (b) montre que $\varphi(\lambda)w=w$ pour tout $\lambda\in C^*$. L'élément $w$ est donc bien invariant par $\Phi$, cqfd.

Soit $M$ le plus petit sous-groupe algébrique de $\mathbf{GL}_V$ (au sens strict du terme, i.e. «défini sur $\mathbf Q_p$») dont le groupe des points à valeurs dans $C$ contienne $\Phi$. Je dirai que $M$ est \textit{le groupe de Mumford-Tate} du module galoisien $V$ (c'est en effet l'analogue $p$-adique de celui défini dans [5]). Le théorème 1 est visiblement équivalent au suivant:

\textbf{Théorème 2.} -- \textit{Le groupe $G_{\mathrm{alg}}$ contient le groupe $M$.}
\margin{8}
